\scanpage{207}
\hfill -- \mspage{52} -- \hfill IX

Comme le foncteur $X\longmapsto\widehat X$ (complétion formelle le long de la fibre spéciale) induit une équivalence entre la catégorie des schémas finis sur $S$, et celle des schémas formels finis sur $S$ (EGA III 4.8), nous allons identifier simplement un schéma fini sur $S$ avec le schéma formel correspondant. Avec cette identification, la définition des ``parties fixes'' nous donne aussitôt pour tout $\nu\geq0$ :
\[
\text{(5.1.4)}\qquad ({}_{\ell^\nu}A^{\unsure{0|o}})^f={}_{\ell^\nu}(\widehat A^{\unsure{0|o}})\quad ,
\]
ce qu'on peut écrire aussi comme une identité de systèmes projectifs
\[
\text{(5.1.5)}\qquad T_\ell(A^{\unsure{0|o}})^f=T_\ell(\widehat A^{\unsure{0|o}})\quad .
\]
On obtient donc un sous-pro-groupe de Barsotti-Tate de la partie fixe de $T_\ell(A^{\unsure{0|o}})$ en prenant $T_\ell(\widehat T)\subset T_\ell(\widehat A^{\unsure{0|o}})$. On définira alors la \uline{partie torique} de $T_\ell(A^{\unsure{0|o}})$ comme
\[
\text{(5.1.6)}\qquad T_\ell(A^{\unsure{0|o}})^t\overset{\mathrm{dfn}}=T_\ell(\widehat T)\subset T_\ell(\widehat A^{\unsure{0|o}})=T_\ell(A^{\unsure{0|o}})^f\quad .
\]
Donc par définition, on aura
\[
\text{(5.1.7)}\qquad T_\ell(A^{\unsure{0|o}})^t=({}_{\ell^\nu}\widehat T)_{\nu\geq0}\quad ,
\]
le $S$-groupe fini et plat ${}_{\ell^\nu}\widehat T$ étant lui-même décrit comme la limite inductive des $S_n$-groupes finis et plats
\[
\text{(5.1.8)}\qquad {}_{\ell^\nu}(T_n),\ n\text{ variable}\quad ,
\]
i.e. son algèbre affine est décrite comme la limite projective des algèbres affines des groupes (5.1.8). Bien entendu, on aura
\[
\text{(5.1.9)}\qquad (T_\ell(A^{\unsure{0|o}})^t)_n=T_\ell(T_n)\quad ,
\]
\margin{370}
\margin{370}
\margin{206}
